\chapter{Вырождение энергий компакттонов и граничный канал перехода}
\label{ch:v6-spectral-transition-discrete-compacton-energy-degeneracy-boundary-overlap-gate}

\section{Различающий вопрос}

Предыдущий гейт показал, что эта/Pfaffian-фаза различает характеры
$+i$ и $-i$, но не даёт положительной скорости. Здесь проверяются четыре
возможности получить предпочтение без вставки $\gamma$:
\begin{enumerate}
 \item разные вещественные энергии или радиусы двух компакттонов;
 \item нелинейное усиление малого перекоса характерных весов;
 \item градиентный сток дефекта $D_\chi$ после проекции $P_3$;
 \item ненулевой граничный матричный элемент между ветвями.
\end{enumerate}

Антикруговая сверка использует
\path{version6_spectral_transition_discrete_compacton_existence_gate},
\path{version6_spectral_transition_discrete_compacton_dynamical_capture_gate},
\path{version6_spectral_transition_compacton_c4_affine_selector_admissibility_gate}
и
\path{version6_spectral_transition_discrete_compacton_c4_boundary_eta_dissipation_gate}.

\section{Объём понятия энергии}

Текущая модель задаёт нелинейный унитарный шаг $F$, но не отдельный
стационарный гамильтонов функционал, значение которого обязано быть
физической массой. Поэтому сравниваются все уже определённые вещественные
локальные данные: норма, центр, квадрат радиуса, обратное участие,
\begin{equation}
 \sum_x p_xq_x,
 \qquad
 \sum_x\|H_{\mathrm{eff},x}\|^2,
 \qquad
 \langle\Psi,K(\Psi)\Psi\rangle,
 \qquad
 \langle\Psi,K(\Psi)^2\Psi\rangle.
 \label{eq:v6-compacton-real-invariants}
\end{equation}
Для
\begin{equation}
 \Psi_+=(v_*,-iv_*),
 \qquad
 \Psi_-=(v_*,+iv_*)
 \label{eq:v6-compacton-conjugate-branches}
\end{equation}
профили комплексно сопряжены. Аудит даёт нулевую разность каждого
функционала~\eqref{eq:v6-compacton-real-invariants}; носитель и радиус
также совпадают точно.

Собственные фазы шага равны $\pm i$, то есть главные квазиэнергетические
углы имеют противоположные знаки и одинаковый модуль $\pi/2$.
Квазиэнергия флоке-системы определена по модулю частоты, поэтому выбор
ветви логарифма сам по себе не является вещественным энергетическим
расщеплением \cite{energy-degeneracy-floquet-review}.

\begin{proposition}[Нет геометрического энергетического предпочтения]
В рамках всех вещественных функционалов, фактически заданных текущим
дискретным родителем, $\Psi_+$ и $\Psi_-$ точно вырождены. Проекция на
компактный носитель также не различает их, поскольку их носители
совпадают.
\end{proposition}

\section{Нелинейность не усиливает характерный перекос}

Просканировано точное фазовое семейство
\begin{equation}
 \Psi_\varphi=(v_*,e^{i\varphi}v_*),
 \qquad 0\leq\varphi<2\pi.
\end{equation}
Для тридцати двух фаз после сорока нелинейных шагов максимальный дрейф
весов $w_{+i},w_{-i}$ равен $8.44\cdot10^{-15}$. Все вещественные
инварианты постоянны с точностью до накопления округления
$1.42\cdot10^{-14}$. Следовательно, равная смесь не является
неустойчивой относительно самофокусировки в пространстве характеров.

Это согласуется с тем, что нелинейные квантовые блуждания могут иметь
солитонные и компактные решения, но их существование само по себе не
гарантирует диссипативного выбора между сопряжёнными флоке-ветвями
\cite{energy-degeneracy-nonlinear-walk}.

\section{Граничное действие также вырождено}

Для скручиваний $\alpha_+=1/4$ и $\alpha_-=3/4$ имеется точное спаривание
\begin{equation}
 \left(n+\frac14\right)^2
 =\left(-n-1+\frac34\right)^2.
 \label{eq:v6-compacton-boundary-spectrum-pairing}
\end{equation}
Поэтому вещественная часть детерминантного действия совпадает. Прямой
скан регуляризованного модуля при пяти значениях радиального параметра
даёт максимальную разность $2.22\cdot10^{-16}$. Эта-инварианты имеют
противоположные знаки, но меняют только фазу.

На аффинном носителе естественная высота $h=-P_3$ имеет одно и то же
значение $-1$ на линиях характеров $i,-1,-i$. Значит, ни граничный
детерминант, ни уже построенная аффинная высота не создают
$E_+\ne E_-$.

\section{Почему скалярная нормальная производная не смешивает ветви}

Пусть $C$ --- выбранный четырёхцикл. Векторы $u_+,u_-$ удовлетворяют
\begin{equation}
 Cu_+=iu_+,
 \qquad
 Cu_-=-iu_-.
\end{equation}
Если граничный оператор $B$ скалярен или $C_4$-эквивариантен, то
$[B,C]=0$, откуда по ортогональности различных характеров
\begin{equation}
 \langle u_+,Bu_-\rangle=0.
 \label{eq:v6-compacton-equivariant-boundary-zero}
\end{equation}
Перебор случайных полиномов $B(C)$ подтверждает равенство с максимальным
остатком $2.27\cdot10^{-15}$. Поэтому предложенный элемент
$\langle+i|\partial_n|-i\rangle$ равен нулю, пока $\partial_n$ является
скалярным граничным оператором.

Однако оператор
\begin{equation}
 T_{+-}=|+i\rangle\langle-i|
\end{equation}
преобразуется как
\begin{equation}
 CT_{+-}C^{-1}=-T_{+-}.
 \label{eq:v6-compacton-minus-character-mixer}
\end{equation}
То есть требуемый смеситель несёт характер $-1$, а не тривиальный
характер. Именно такая третья линия уже существует внутри аффинного
$P_3\mathbb C^4$ наряду с линиями $\pm i$.

\begin{theorem}[Найден минимальный посредник]
Скалярная $C_4$-граница не может породить переход
$|-i\rangle\to|+i\rangle$. Минимальный эквивариантный посредник должен
лежать в характере $-1$. Такой характер уже содержится в аффинном
триплете проекта, поэтому новый внутренний тип представления вводить не
нужно.
\end{theorem}

Это ещё не диссипация. Конечная динамическая мода характера $-1$ даст
когерентные осцилляции. Необратимый rate может появиться лишь если эта
мода имеет непрерывный спектр, исходящий канал или после её исключения
возникает мнимая часть собственной энергии.

\section{Что даёт проецированный градиент}

Для
\begin{equation}
 D_\chi=4w_+(1-w_+)
\end{equation}
естественный симплексный градиентный закон можно записать без отдельного
безразмерного коэффициента:
\begin{equation}
 \frac{dw_+}{ds}
 =4w_+(1-w_+)(2w_+-1).
 \label{eq:v6-compacton-dchi-gradient}
\end{equation}
Точка $w_+=1/2$ остаётся неподвижной, но любой заданный перекос выбирает
одну из чистых ветвей. Численно $0.49$ и $0.51$ уходят к разным концам.

Однако~\eqref{eq:v6-compacton-dchi-gradient} является новым
градиентным родительским законом. Переименование времени в $s$ скрывает,
но не выводит физическую метрику и перевод между $s$ и наблюдаемым
временем. Поэтому $P_3D_\chi$ фиксирует направление стока, но не его
физическую скорость.

\begin{nogocriterion}[Скрытый коэффициент трения]
Нельзя считать коэффициент выведенным только потому, что он поглощён в
безразмерное время градиентного потока. Требуется получить метрику потока
или спектральную плотность исключаемой моды из уже заданного действия.
\end{nogocriterion}

\section{Вердикт}

Чистейшая попытка получить предпочтение из разных энергий закрыта:
компакттоны $\pm i$ точно вырождены по всем имеющимся вещественным
объёмным и граничным данным. Текущая нелинейность сохраняет их веса, а
скалярный граничный оператор не имеет внедиагонального матричного
элемента. Градиент $D_\chi$ выбирает ветвь только после объявления нового
диссипативного закона.

Но найдено конкретное продолжение без произвольного оператора перехода:
требуемый смеситель имеет характер $-1$, и эта линия уже присутствует в
аффинном $P_3$-триплете. Следующий гейт
\path{version6_spectral_transition_discrete_compacton_character_resolved_radiation_form_factor_gate}
должен проверить её непосредственно в полном флоке-якобиане и вычислить
форм-фактор в уже существующий пространственный радиационный континуум.

\begin{protocol}
\begin{enumerate}
 \item разность всех заданных вещественных объёмных инвариантов:
 \textbf{ноль};
 \item радиусы и носители $\Psi_\pm$: \textbf{совпадают};
 \item абсолютные квазиэнергетические углы: \textbf{совпадают};
 \item физический стационарный гамильтониан текущей модели:
 \textbf{не задан};
 \item нелинейное усиление перекоса $w_+-w_-$: \textbf{нет};
 \item вещественные граничные действия: \textbf{вырождены};
 \item скалярный граничный переход $\pm i$: \textbf{запрещён};
 \item характер минимального смесителя: \textbf{$-1$};
 \item линия $-1$ в существующем $P_3$-триплете: \textbf{есть};
 \item её динамический профиль и открытый спектр: \textbf{не выведены};
 \item градиент $D_\chi$ выбирает ветвь: \textbf{да, условно};
 \item физическая метрика и скорость градиента: \textbf{не выведены};
 \item $\gamma$ устранён из физической динамики: \textbf{нет};
 \item новый структурный маршрут без ручного jump-оператора:
 \textbf{найден}.
\end{enumerate}
\end{protocol}

Машинный сертификат сохранён в
\path{s2t_v6_spectral_transition_discrete_compacton_energy_degeneracy_boundary_overlap_gate_results.json}.

\begin{thebibliography}{9}
\bibitem{energy-degeneracy-floquet-review}
M.~S.~Rudner and N.~H.~Lindner,
\emph{Floquet topological insulators: from band structure engineering to novel non-equilibrium quantum phenomena},
Nature Reviews Physics 2 (2020), 229--244,
\path{arXiv:1909.02008}.

\bibitem{energy-degeneracy-nonlinear-walk}
C.-W.~Lee, P.~Kurzyński, and H.~Nha,
\emph{Quantum walk as a simulator of nonlinear dynamics: Nonlinear Dirac equation and solitons},
Physical Review A 93 (2016), 042325,
\path{arXiv:1512.08358}.
\end{thebibliography}